\section{1988--1998: Becoming an Institution}

\subsection{Wigratzbad: a seminary before a society}

The papal rescript of 18 October 1988 authorized the erection of a
seminary at Wigratzbad, in the township of Opfenbach in the diocese of
Augsburg, ``with the prior consent of the diocesan bishop.'' The
Fraternity's own history of the house records that the consent came
first: ``As early as August, 1988, Bishop Josef Stimpfle of Augsburg
announced that he would accept a canonical establishment in Wigratzbad.
In November thirty-one seminarians commenced their studies.''\endnote{``The
Seminary of Wigratzbad (Germany)'' (fssp.org, ``Our seminaries''), read
2026-07-25. A self-published institutional history, undated on the page
and stale in one respect noted below. Bishop Stimpfle's own act was not
located; the rescript of 18 October 1988 establishes that the bishop's
prior consent was required, not that it was given in August, and this
account carries the August date on the Fraternity's authority alone.}

The site was not chosen at random. Wigratzbad was already a Marian
pilgrimage centre associated with Antonie Rädler and, in the Fraternity's
telling, a place where the idea of an international seminary had been
discussed years before 1988. Until autumn 2000 the seminary was housed in
the pilgrimage centre itself; a purpose-built seminary opened at the
start of the 2000 term and was solemnly blessed by Cardinal Darío
Castrillón Hoyos on 2 December 2000. The Fraternity records that from the
beginning the studies were divided into French and German language
groups, with seminarians of other languages choosing one of the
two.\endnote{Same page. The date of the blessing---2 December 2000---is
worth holding: it falls five months after the Roman intervention in the
Fraternity's government described below, and the prelate who blessed the
new building is the prelate who had made that intervention.}

Two things follow from the shape of that first decade. The Fraternity was
an international body before it was a large one, because a single house
of formation served several language groups; and its centre of gravity
was formation rather than parochial ministry, which is what its own
constitutions later say its principal object is.

\subsection{Elmhurst, then Denton}

The second seminary has a two-stage history that the Fraternity's own
account preserves and that is easily lost in summary. In 1994 Bishop
James C.\ Timlin of Scranton, Pennsylvania approved the establishment of
a seminary in a former retreat centre at Elmhurst, and that autumn the
first seminarians began studies at Our Lady of Guadalupe Seminary. By
1998 the numbers had outgrown the site; Father Arnaud Devillers, then
North American district superior, announced that ``the new seminary will
be located in the Diocese of Lincoln, Nebraska where we have been
welcomed by Bishop Fabian W.\ Bruskewitz.'' Ground was broken in October
1998; Bishop Bruskewitz blessed the cornerstone on 16 October 1999; and
in autumn 2000 the new house at Denton, Nebraska received nearly fifty
seminarians under its first rector, Father James Jackson. The seminary
chapel of Saints Peter and Paul was consecrated by Bishop Bruskewitz on 3
March 2010 in the presence of Cardinal William Levada, then Prefect of the
Congregation for the Doctrine of the Faith and President of the
Commission \work{Ecclesia Dei}.\endnote{``The Seminary of Denton (USA)''
(fssp.org, ``Our seminaries''), read 2026-07-25, from which the quoted
words of Father Devillers and Bishop Bruskewitz are taken. Self-published
institutional history; the diocesan instruments of Scranton and Lincoln
were not located. Note that the page is stale---it describes the
building's history ``nearly a decade ago'' and names a rector who has
since been succeeded---which is a reason to use it for dated events and
not for present facts. A widely circulated agency report of January 2026
places the Denton seminary in Texas; it is in Nebraska.}

The move from Scranton to Lincoln is worth a sentence of analysis. A
society of apostolic life cannot open a house anywhere it likes: canon
733 §\,1 requires the previous written consent of the diocesan bishop.
The Fraternity's two seminaries therefore exist because three named
bishops consented---Stimpfle at Augsburg, Timlin at Scranton and then
Bruskewitz at Lincoln---and its whole apostolic network exists on the
same terms. A
body of pontifical right is not thereby free of bishops; it is a body
whose houses each rest on an episcopal act.

\subsection{The rite of ordination, and a rescript of 1996}

One institutional fact of this period is easy to overlook and is
canonically important. A society that forms its own clergy must present
them for orders, and the erection decree had conceded the 1962 books,
which include the Pontifical. The canonical literature records a further
rescript of the Pontifical Commission of 24 March 1996, Prot.\ N.~64/96,
by which the Fraternity could celebrate the rite of tonsure and of
ordination according to the \work{Pontificale Romanum} of
1961/62.\endnote{Pietras, art.\ cit., 176, reporting the rescript with its
protocol number. Bounded negative: the rescript itself was not located,
and no text of it is published by the Fraternity or by the Holy See. The
fact is reported here on the authority of a peer-reviewed canonical
article whose author is a priest of the Fraternity; the qualification is
stated because the reader should know that the only witness to this act is
an interested one, even if a scholarly one.} The point is not
antiquarian. Twenty-five years later the Congregation for Divine Worship
would answer that the diocesan bishop may not permit the pre-conciliar
Pontifical at all. Whatever the Fraternity holds in that matter, it holds
by a distinct title.

\subsection{Growth, and what the numbers are}

The Fraternity publishes annual statistics as of 1 November each year.
They are self-reported; no independent census exists and none is
attempted here. Taken as a self-report, they show a body that grew
by about a third between 2017 and 2024 and whose seminary population
never collapsed:

\begin{documentinventory}
\inventoryrow{1 November 2017}{Fraternity statistics page}{437 members,
of whom 293 incardinated; 287 priests; 21 deacons; 129 non-deacon
seminarians; average age 38.}
\inventoryrow{1 November 2019}{Fraternity statistics page}{482 members,
317 incardinated; 320 priests; 17 deacons; 145 non-deacon seminarians;
average age 38.}
\inventoryrow{1 November 2021}{Fraternity statistics page}{526 members,
336 incardinated; 341 priests; 17 deacons; 168 non-deacon seminarians;
average age 38.}
\inventoryrow{1 November 2023}{Fraternity statistics page}{569 members,
364 incardinated; 368 priests; 22 deacons; 179 non-deacon seminarians;
average age 39.}
\inventoryrow{1 November 2024}{Fraternity statistics page}{583 members,
373 incardinated; 386 priests; 15 deacons; 182 non-deacon seminarians;
average age 39.}
\inventoryrow{1 November 2025}{Fraternity statistics page}{579 members,
394 incardinated; 387 priests (365 incardinated, 16 incorporated
\latin{ad annum}, 3 postulants, 3 associated); 30 deacons; 162 non-deacon
seminarians; average age 39; 16 deceased members.}
\end{documentinventory}

Three cautions travel with that table.\endnote{``Statistics update'' for
1 November 2017, 2019, 2021, 2023, 2024 and 2025, and the standing
``Figures'' page (fssp.org), all read 2026-07-25. Self-reported;
categories are the Fraternity's own. The 2025 page also reports priestly
ordinations over twelve years, from a low of 11 (2021 and 2025) to a high
of 19 (2017), averaging 14 a year, and gives the lay Confraternity of
Saint Peter 8{,}776 members. The ``deceased members'' figure is
cumulative since foundation, not annual.} First, the categories are the
Fraternity's own and are not interchangeable: ``member'' includes
seminarians, ``incardinated'' does not, and ``priests'' includes men
incorporated \latin{ad annum}, postulants and associates who are not
incardinated. Second, the 2025 figure is the first fall in the series
reported here, from 583 to 579, driven by a drop of twenty in non-deacon
seminarians; a single year is not a trend and this account draws no
conclusion from it. Third, membership is not the same as apostolic
extent: the canonical literature reports 151 dioceses, 277 apostolates and
48 personal parishes at 2024, drawn from the Fraternity's own
\work{Ordo administrativus}.\endnote{Pietras, art.\ cit., 166 n.~1, citing
the 2024 \work{Ordo administrativus}. That internal directory was not
reached for this account, and the figures are reported here as the
article reports them.}

\subsection{Government, and one unresolved date}

The Fraternity is governed by a superior general with a general council,
elected by a general chapter, and is divided territorially. Its published
succession of superiors general runs Josef Bisig (1988--2000), Arnaud
Devillers (2000--2006), John Berg (2006--2018), Andrzej Komorowski
(2018--2024), and John Berg again from 2024, elected on 9 July 2024 by
the seventh ordinary general chapter, at which thirty-two capitulants were
present.\endnote{``Election of the Superior General of the Fraternity''
(11 July 2024) and ``Closing of the General Chapter and nominations''
(18 July 2024), both fssp.org, read 2026-07-25; ``A new Superior General
for the Fraternity'' (10 July 2018) for the Komorowski election, which
that communiqué describes as making him ``the 4th Superior General.'' The
1988--2000 and 2000--2006 terms are established by the documents discussed
in the sections on 1999 and 2000 below.}

The territorial history contains a discrepancy that this account does not
resolve. The Fraternity's communiqué of 27 May 2021 states that
``following the decision of the 2018 General Chapter, the North American
District (USA and Canada) will also be erected into a North American
Province on July 1''; the canonical article of 2026 gives ``the American
Province, established in 2002.'' A North American \emph{district} plainly
existed earlier still, since the Fraternity's own seminary history calls
Father Devillers ``then North American District Superior'' in 1998. Both
statements cannot describe the same act, neither writer supplies the
erecting instrument, and this account does not choose between them. What
is established is that a province of that name existed at the cutoff and
that the Fraternity dates its erection to 1 July 2021.\endnote{``Appointments and new structures within the FSSP''
(fssp.org, 27 May 2021), read 2026-07-25, which also records the erection
of the Southern Cross Region (Australia and New Zealand) into a district
from 1 July 2021; Pietras, art.\ cit., 166 n.~1; and ``The Seminary of
Denton (USA)'' (fssp.org) for the description of Father Devillers as
North American district superior in 1998. Neither writer supplies an
erecting instrument, and none was located.}

At the cutoff the Fraternity's published organization chart shows a
superior general (John Berg) with three assistants and two counsellors, a
North American Province, a French Province, German-speaking and Oceania
districts, and the two seminaries.\endnote{``Organization chart''
(fssp.org, ``Presentation''), read 2026-07-25. Bounded negative: no
instrument erecting a French Province was located, and the chart is the
only witness to its existence reached for this account. The chart is a
current-state page with no date on its face, which is a further reason to
treat it as a snapshot of 25 July 2026 and nothing more.}
